\documentclass[12pt, a4paper, twoside]{article}
\usepackage[a4paper, margin=1in]{geometry}
\usepackage{xeCJK}
\usepackage{graphicx}
\usepackage{subcaption}
\usepackage{caption}
\usepackage{amsmath}
\usepackage{verbatim}
\usepackage{listings}
\usepackage{xcolor}
\usepackage{url}
\usepackage{float}
\usepackage{booktabs}
\usepackage{array}
\usepackage{amssymb}
\usepackage{multirow}
\usepackage{tikz}
\usepackage{enumitem}
\usepackage{cite}

\usepackage{indentfirst} %首行縮排
\setlength{\parindent}{2em}

% 設定中英文字型
\setCJKmainfont{Noto Sans CJK TC}  % 中文：正黑體
%\setmainfont{Arial}        % 英文字體
% \setCJKmonofont{Microsoft JhengHei}  % 中文等寬字型
% \setmonofont{Courier New}            % 英文等寬字型

\title{AAAA}
\author{Your Name} 
\date{\today}

\renewcommand{\baselinestretch}{1.0}

\begin{document}
\maketitle
%\raggedright
%\noindent % 避免段落縮排

\begin{abstract}


\end{abstract}

\section{Introduction}
I had browsed through existing works in recent 3 years in the field of Intelligent Vehicle and summarized the topic and the trend of the research. \cite{lei2022latency, hu2022where2comm, xia2024one, phcp2025progressive, wang2026communication, wei2025cooperative, wang2025latency, poibrenski2025robust, gao2025spatiotemporal, wang2025towards, huang2025vehicle, guo2025when}

There are several main problems in the field of Intelligent Vehicle, including bandwidth \& communication bottlenecks, latency \& temporal asynchrony, heterogeneity \& cross-model compatibility, infrastructure-assisted efficiency (V2I), robustness, security \& reliability, and sim-to-real gap \& empirical deployment constraints. 

Then I focus on the problem of latency \& temporal asynchrony in cooperative perception.
These two challenges are critical to spatial-temporal data consistency.
Latency refers to the delay form sensation to fusion and decision-making, which can be caused by network transmission, data processing, and algorithmic computation.
This latency come from bottlenecks in communication, computation, and synchronization, which can lead to outdated or incomplete information being used for decision-making.
Temporal asynchrony refers to the misalignment of data timestamps from different sources,including RSUs and CAVs, which can lead to inconsistencies in the perception results. 

Usually, those data fused trivially by addition or concatention, which can lead to spatial-temporal misalignment and inaccurate perception results.
Moving objects are predicted to be in wrong historical positions, instead of their current positions, leading to phantom detections, tracking fragmentation, and dangerous miscalculations in downstream motion planning.

\section{}

% 設定 BibTeX 參考文獻樣式與引用資料庫檔名 
\bibliographystyle{IEEEtran} 
% 可依投稿期刊更換為 plain, unsrt, ACM-Reference-Format 等 
\bibliography{ref} % 載入 references.bib
\end{document}